\section{Agent-Based Model Equations}
\label{sec:abm_equations}

The crop--soil system is modeled as a set of $N = 130$ coupled discrete-time dynamical systems, building upon the ABM framework introduced by L\'{o}pez-Jim\'{e}nez et al.~\cite{Lopezjimenez2024}. Each agent $n \in \{1, \ldots, N\}$ maintains four state variables that evolve at a daily time step $k$. These states capture the essential dynamics of the soil--plant--atmosphere continuum: water availability, thermal accumulation, crop maturation, and biomass production.

To rigorously enforce mass conservation and physical realism, the hydrological dynamics are decoupled into a surface layer and a subsurface root-zone layer, bridged by an infiltration term.

\subsection{State Equations}
\label{sec:state_equations}

The four state equations governing the evolution of agent $n$ are:
\begin{align}
    x_1^n(k+1) &= x_1^n(k) + I^n(k) + \varphi_{4,\text{in}}^n(k) - \varphi_1^n(k) - \varphi_3^n(k) - \varphi_{4,\text{out}}^n(k) \label{eq:x1} \\
    x_2^n(k+1) &= x_2^n(k) + h_1^n(k) \label{eq:x2} \\
    x_3^n(k+1) &= x_3^n(k) + \theta_{11}\bigl(1 - h_2(k)\bigr) + \theta_{12}\bigl(1 - h_3^n(k)\bigr) \label{eq:x3} \\
    x_4^n(k+1) &= x_4^n(k) + \theta_{13} \cdot h_3^n(k) \cdot h_6^n(k) \cdot h_7(k) \cdot g(x_2^n) \cdot R_s(k) \label{eq:x4}
\end{align}

where $x_1$ is the volumetric soil water content in the root zone (mm), $x_2$ is the cumulative thermal time ($^{\circ}$C$\cdot$day), $x_3$ is the cumulative maturation index incorporating heat and drought stress effects, and $x_4$ is the aboveground dry biomass (g/m$^2$). The environmental inputs are daily precipitation $P(k)$ (mm/day) and solar radiation $R_s(k)$ (MJ/m$^2$/day). The control input $u^n(k)$ is the irrigation applied to agent $n$ on day $k$ (mm/day). 

\subsection{Surface and Subsurface Water Fluxes}
\label{sec:soil_fluxes}

\textbf{Surface Layer and Infiltration:} 
The total available surface water $W_{\text{surf}}$ consists of precipitation, applied irrigation, and incoming surface runoff from higher-elevation neighbors ($\varphi_{2,\text{in}}$). 
\begin{equation}
    W_{\text{surf}}^n(k) = P(k) + u^n(k) + \varphi_{2,\text{in}}^n(k)
    \label{eq:wsurf}
\end{equation}

The incoming surface and subsurface flows for an agent $n$ are the sum of the respective outflows from all uphill neighbors $m \in \mathcal{U}(n)$, divided equally among the neighbor's downhill receivers $N_r(m)$:
\begin{align}
    \varphi_{2,\text{in}}^n(k) &= \sum_{m \in \mathcal{U}(n)} \frac{\varphi_2^m(k)}{N_r(m)} \\
    \varphi_{4,\text{in}}^n(k) &= \sum_{m \in \mathcal{U}(n)} \frac{\varphi_{4,\text{out}}^m(k)}{N_r(m)}
\end{align}

Surface runoff ($\varphi_2$) is computed using the SCS curve number method, activated when surface water exceeds the initial abstraction threshold $\theta_3$:
\begin{equation}
    \varphi_2^n = 
    \begin{cases}
        \dfrac{(W_{\text{surf}}^n - \theta_3)^2}{W_{\text{surf}}^n + 4\theta_3}, & W_{\text{surf}}^n > \theta_3 \\[8pt]
        0, & W_{\text{surf}}^n \leq \theta_3
    \end{cases}
    \label{eq:phi2}
\end{equation}
Water that does not run off the surface infiltrates into the root zone:
\begin{equation}
    I^n(k) = W_{\text{surf}}^n(k) - \varphi_2^n(k)
    \label{eq:infiltration}
\end{equation}

\textbf{Crop transpiration} ($\varphi_1$) is the water extracted by plant roots to meet evaporative demand, bounded by the reference evapotranspiration ($\text{ET}_0$):
\begin{equation}
    \varphi_1^n = \min\!\Big(\max\!\big(\theta_1(x_1^n - \theta_2 \theta_5), 0\big),\; \text{ET}_0\Big)
    \label{eq:phi1}
\end{equation}
where $\theta_1$ is the water uptake coefficient, $\theta_2$ is the volumetric wilting point, and $\theta_5$ is the root zone depth (mm). 

\textbf{Subsurface Excess Partitioning:}
To calculate deep drainage and lateral subsurface flow, an intermediate soil moisture state ($x_{1,\text{temp}}$) is computed after infiltration, transpiration, and incoming lateral subsurface flow ($\varphi_{4,\text{in}}$) are resolved:
\begin{equation}
    x_{1,\text{temp}}^n(k) = x_1^n(k) + I^n(k) + \varphi_{4,\text{in}}^n(k) - \varphi_1^n(k)
\end{equation}
If this intermediate moisture exceeds the field capacity ($\theta_6 \theta_5$), the subsurface excess $E_{\text{sub}}^n$ is forced out of the root zone:
\begin{equation}
    E_{\text{sub}}^n(k) = \max\!\big(x_{1,\text{temp}}^n(k) - \theta_6 \theta_5,\; 0\big)
\end{equation}
This excess is strictly partitioned into vertical \textbf{deep drainage} ($\varphi_3$) and horizontal \textbf{lateral outflow} ($\varphi_{4,\text{out}}$) using the drainage coefficient $\theta_4$:
\begin{align}
    \varphi_3^n(k) &= \theta_4 \cdot E_{\text{sub}}^n(k) \label{eq:phi3} \\
    \varphi_{4,\text{out}}^n(k) &= (1 - \theta_4) \cdot E_{\text{sub}}^n(k) \label{eq:phi4out}
\end{align}

To resolve the spatial dependency of lateral flows within a single time step $k$, the state equations are solved sequentially following a topological sort of the directed elevation graph, executing from the highest elevation agents down to the sink agents.

\subsection{Stress and Growth Functions}
\label{sec:stress_growth}

\textbf{Thermal time accumulation} ($h_1$) drives phenological development for temperatures above the base threshold $\theta_7$:
\begin{equation}
    h_1^n = \max\!\bigl(T_{\text{mean}} - \theta_7,\; 0\bigr)
    \label{eq:h1}
\end{equation}

\textbf{Heat stress} ($h_2$) reduces maturation efficiency linearly between an optimal temperature ceiling ($\theta_9$) and an extreme heat threshold ($\theta_{10}$):
\begin{equation}
    h_2(k) = \text{clip}\!\left( \frac{\theta_{10} - T_{\max}(k)}{\theta_{10} - \theta_9},\; 0,\; 1 \right)
\end{equation}
\textbf{Low temperature stress} ($h_7$) prevents biomass accumulation when daily temperature drops below the base threshold:
\begin{equation}
    h_7(k) = \begin{cases} 
      1, & T_{\text{mean}}(k) > \theta_7 \\ 
      0, & T_{\text{mean}}(k) \leq \theta_7 
   \end{cases}
\end{equation}

\textbf{Drought stress} ($h_3$) quantifies the reduction in crop performance when actual transpiration falls short of atmospheric demand:
\begin{equation}
    h_3^n = 1 - \theta_{14} \cdot \max\!\left(1 - \frac{\varphi_1^n}{\text{ET}_0},\; 0\right)
    \label{eq:h3}
\end{equation}
where $\theta_{14}$ controls the crop's sensitivity to drought. 

\textbf{Waterlogging stress} ($h_6$) penalizes biomass growth when soil water exceeds field capacity, simulating root-zone oxygen deprivation:
\begin{equation}
    h_6^n = \text{clip}\!\left(1 - \frac{x_1^n - \theta_6 \theta_5}{\theta_6 \theta_5},\; 0,\; 1\right)
    \label{eq:h6}
\end{equation}

\textbf{Growth function} ($g$) governs the fraction of intercepted radiation contributing to biomass via a two-branch sigmoid model driven by cumulative thermal time ($x_2$):
\begin{equation}
    g(x_2^n) = 
    \begin{cases}
        \dfrac{\theta_{19}}{1 + \exp\!\bigl(-0.01(x_2^n - \theta_{20})\bigr)}, & x_2^n \leq \dfrac{\theta_{18}}{2} \\[14pt]
        \dfrac{\theta_{19}}{1 + \exp\!\bigl(0.01(x_2^n + x_3^n - \theta_{18})\bigr)}, & x_2^n > \dfrac{\theta_{18}}{2}
    \end{cases}
    \label{eq:growth}
\end{equation}


\section{Model Corrections and Extensions}
\label{sec:corrections}

While the baseline framework~\cite{Lopezjimenez2024} provides a robust topographical architecture, implementation for closed-loop control revealed necessary mathematical corrections to enforce strict mass-balance and to ensure controller cost-functions were properly coupled to agronomic penalties. Three primary modifications were integrated into the environment.

\subsection{Mass Conservation and Hydrological Partitioning}

The original formulation in~\cite{Lopezjimenez2024} computes an outflow term ($\phi_{out}$) that lumps surface runoff and subsurface excess together. While the runoff portion is properly subtracted from the sender's state equation, the subsurface excess portion is not. Consequently, a sending agent passes its subsurface excess to downhill neighbors but never loses that volume from its own root zone. This generates an unbounded artificial creation of water ("phantom water"), leading to exaggerated flooding in lower elevations during sustained rainfall. 

Furthermore, because vertical deep drainage ($\varphi_3$) and the subsurface portion of the lateral outflow both independently drew from the total excess pool in the original model, the fluxes mathematically overlapped. This allowed the system to subtract more than 100\% of the available excess water from the root zone.

The proposed decoupled two-layer architecture resolves these flaws. By explicitly bridging the surface and subsurface via infiltration ($I$), and by strictly partitioning the subsurface excess ($E_{\text{sub}}$) using the drainage coefficient ($\theta_4$), horizontal and vertical fluxes are perfectly conserved. The explicit subtraction of $\varphi_{4,\text{out}}^n$ guarantees strict mass conservation.

\subsection{Drought Stress on Biomass}

In the original model, the drought stress function $h_3$ appears only in the maturation equation ($x_3$, Equation~1c of~\cite{Lopezjimenez2024}), where it accelerates crop senescence under water deficit. However, $h_3$ is absent from the biomass equation ($x_4$, Equation~1d). This means that a completely unirrigated crop would accumulate the same daily biomass increment as a fully irrigated one, provided all other conditions were equal. The only effect of drought would be to advance the onset of senescence, which reduces biomass accumulation only indirectly and with considerable delay.

This omission was identified during the initial ABM validation, when the no-irrigation baseline produced biomass values that were nearly indistinguishable from irrigated scenarios for wheat. The corrected formulation multiplies the radiation use efficiency term in Equation~\eqref{eq:x4} by $h_3$, so that drought directly reduces the daily biomass increment. This modification follows the approach of the FAO AquaCrop model, in which water stress reduces radiation use efficiency proportionally to the ratio of actual to potential transpiration. With $\theta_{14} = 0.8$ for rice and 0.6 for tobacco, maximum drought reduces the biomass increment by up to 80\% and 60\%, making irrigation decisions directly consequential for yield outcomes.

This correction is particularly important for the controller comparison objective of this thesis. Without it, the ABM would produce only marginal differences between irrigated and unirrigated outcomes, rendering any comparison between MPC, RL, and baseline controllers meaningless. The original crop used in the baseline paper was wheat in Colombia, which is less water-sensitive (lower $\theta_{14}$); it is possible that the omission went undetected in that context because the irrigation response was already small.


\subsection{Surface Runoff Decoupling and Redistribution}
\label{sec:runoff}

In the original configuration~\cite{Lopezjimenez2024}, surface runoff ($\varphi_2$) is mathematically aggregated with subsurface flow and routed directly into the soil matrix of adjacent agents. This simplification bypasses surface pooling dynamics, injecting rapid overland flow directly into the root zone of downhill neighbors. 

To more accurately capture the bowl-shaped topography of the field, surface runoff is explicitly decoupled from subsurface flow in this work. Runoff generated by an agent is routed to lower-elevation neighbors where it arrives as surface inflow ($\varphi_{2,\text{in}}^n$). This inflow is added to the neighbor's available surface water ($W_{\text{surf}}$) and must undergo infiltration before entering the root zone.

\begin{table}[H]
\centering
\caption{Runoff redistribution comparison, no irrigation, wet scenario (2024).}
\label{tab:runoff}
\begin{tabular}{lccc}
\toprule
\textbf{Mode} & \textbf{Yield (kg/ha)} & \textbf{$x_1$ sinks (mm)} & \textbf{$x_1$ hills (mm)} \\
\midrule
None (original) & 2135 & 71.2 & 71.2 \\
Simple (A) & 2273 & 74.8 & 72.0 \\
Cascade (B) & 2278 & 79.5 & 72.0 \\
\bottomrule
\end{tabular}
\end{table}

A cascading alternative (top-to-bottom SCS re-application within a single time step) was also tested. As demonstrated in Table \ref{tab:runoff}, the simple single-step redistribution captures the primary agronomic effect (a 6.5\% yield improvement due to increased water pooling and retention) with negligible difference from the cascade approach (0.2\%). Therefore, the simple redistribution was adopted to minimize computational overhead during MPC optimization.

\subsection{Biomass Calibration}

The radiation use efficiency parameter $\theta_{13}$ was calibrated so that $x_4$ reliably represents g/m$^2$ of aboveground dry matter. 

Yield: $Y = x_4 \times \text{HI} \times 10$ (kg/ha). 


Rice: validated against $\sim$2000~kg/ha (farmer reports) and 3500--4000~kg/ha (optimized)~\cite{jalali2021}. Initial conditions: rice $x_{2,\text{init}} = 210$, $x_{4,\text{init}} = 60$; tobacco $x_{2,\text{init}} = 0$, $x_{4,\text{init}} = 8$.

\section{Soil and Crop Parameterization}
\label{sec:soil_crop_params}

\subsection{Soil Parameters}
\label{soil_param}
The soil at the study site is classified as silty loam, based on regional soil surveys and FAO soil classification maps for Gilan Province. The soil hydraulic properties used in this study follow FAO Irrigation and Drainage Paper No.~56 recommendations for this soil type. 

\begin{table}[htbp]
\centering
\caption{Soil hydraulic parameters for silty loam (common to all crops)}
\label{tab:soil}
\begin{tabular}{cccc}
\toprule
\textbf{Parameter} & \textbf{Symbol} & \textbf{Value} & \textbf{Source}\\
\midrule
Water uptake coefficient & $\theta_1$ & 0.096 & \cite{Lopezjimenez2024} \\
Wilting point (volumetric fraction) & $\theta_2$ & 0.15 & \cite{allen1998} \\
SCS threshold (mm) & $\theta_3$ & 10.0 & \cite{Lopezjimenez2024} \\
Drainage coefficient & $\theta_4$ & 0.05 & \cite{Lopezjimenez2024} \\
Field capacity (volumetric fraction) & $\theta_6$ & 0.35 & \cite{allen1998} \\
\bottomrule
\end{tabular}
\end{table}

\subsection{Crop Parameters}
\label{crop_param}

Two crops were selected for the simulation environment: rice (Hashemi cultivar) and tobacco (\emph{Nicotiana tabacum}). Both hold significant commercial value in Gilan Province~\cite{surfiran2025, iranica_tobacco}. Their contrasting agronomic characteristics (high water demand, shallow roots, and high drought sensitivity for rice, versus moderate water demand, deep roots, and relative drought tolerance for tobacco) provide a robust basis for evaluating the control algorithms under differing physical constraints.

\subsubsection{Rice (Hashemi Cultivar)}

The base temperature ($T_{\text{base}}$ or $\theta_7$) is set to $10^{\circ}$C, following standard Growing Degree Day (GDD) conventions~\cite{mcmaster1997, laohasiriwong2023} and the baseline ABM~\cite{Lopezjimenez2024}. The maturity threshold ($\theta_{18}$) was calibrated to 1250~$^{\circ}$C$\cdot$day based on the 25-year dataset and validated against local field measurements, which report 1216--1352~GDD for the Hashemi cultivar in Gilan~\cite{sadidi2021}. 

Transplanting is simulated on May 20 (DOY 141), aligning with local practices when the 7-day rolling average $T_{\text{mean}} \geq 18^{\circ}$C and $T_{\min} \geq 12^{\circ}$C~\cite{paredes2025, faghani2011}. The harvest is set for August 20 (DOY 233). Extending the season beyond this date sharply increases the risk of extreme WMO R20mm rainfall events ($\geq$20~mm/day), which can cause crop lodging and severe harvest disruption~\cite{sarma2018, nayak2011}. 

This yields a 93-day field season. As established in Section~\ref{sec:corrections}, the 20-day prior nursery period contributes an initial thermal accumulation of $x_{2,\text{init}} = 210$~GDD and an initial biomass of $x_{4,\text{init}} = 60$~g/m$^2$. Other key parameters include a crop coefficient ($K_c$) of 1.15, a shallow root depth ($\theta_5$) of 400~mm, and a high drought sensitivity ($\theta_{14}$) of 0.80.

\subsubsection{Tobacco (\emph{Nicotiana tabacum})}

The base temperature is similarly set to $10^{\circ}$C, which falls within the standard 10--13$^{\circ}$C restriction range for tobacco~\cite{li2021}. The maturity threshold ($\theta_{18}$) is calibrated to 1200~GDD. This duration corresponds to a 104-day field season, consistent with standard FAO timelines (90--120 days)~\cite{fao_tobacco} and general agronomic data (100--130 days)~\cite{britannica1999}. As tobacco seedlings are typically delivered directly from nurseries with unknown prior conditions, the initial states are set to $x_{2,\text{init}} = 0$ and a minimal $x_{4,\text{init}} = 8$~g/m$^2$.


Transplanting occurs on May 25 (DOY 146), with harvest on September 5 (DOY 249). Tobacco features a significantly deeper root zone ($\theta_5 = 700$~mm) compared to rice, allowing it to access a larger reservoir of deep soil moisture. It also exhibits a lower crop coefficient ($K_c = 0.90$) and lower drought sensitivity ($\theta_{14} = 0.60$), making it more resilient to constrained irrigation schedules. 

\subsubsection{Parameter Summary}

\begin{table}[H]
\centering
\caption{Soil and crop-specific parameters used in the ABM simulation}
\label{tab:params}
\small
\begin{tabular}{llccl}
\toprule
\textbf{Parameter} & \textbf{Symbol} & \textbf{Rice} & \textbf{Tobacco} & \textbf{Source} \\
\midrule
Root depth & $\theta_5$ & 400 mm & 700 mm & \cite{allen1998}, \cite{fao_tobacco} \\
Base temperature & $\theta_7$ & 10$^{\circ}$C & 10$^{\circ}$C & \cite{laohasiriwong2023}, \cite{li2021} \\
Heat stress onset & $\theta_9$ & 35$^{\circ}$C & 30$^{\circ}$C & \cite{fao_tobacco} \\
Extreme heat & $\theta_{10}$ & 42$^{\circ}$C & 38$^{\circ}$C & -- \\
Heat stress maturation rate & $\theta_{11}$ & 0.0030 & 0.0025 & ? \\
Drought stress maturation rate & $\theta_{12}$ & 0.0030 & 0.0025 & ? \\
RUE & $\theta_{13}$ & 0.65 & 0.30 & \cite{patel2013}, \cite{ferrara2002} \\
Drought sensitivity & $\theta_{14}$ & 0.80 & 0.60 & \cite{Lopezjimenez2024} \\
Maturity GDD & $\theta_{18}$ & 1250 & 1200 & \cite{sadidi2021}, calib. \\
Max interception & $\theta_{19}$ & 0.95 & 0.90 & \cite{Lopezjimenez2024} \\
50\% intercept. GDD & $\theta_{20}$ & 417 & 450 & \cite{zhang2017} \\
Crop coefficient & $K_c$ & 1.15 & 0.90 & \cite{allen1998}, \cite{fao_tobacco} \\
Depletion fraction & $p$ & 0.20 & 0.50 & \cite{allen1998}, \cite{fao_tobacco} \\
Harvest index & HI & 0.42 & 0.55 & \cite{jalali2021}, \cite{fao_tobacco} \\
Nursery GDD & $x_{2,\text{init}}$ & 210 & 0 & Calib. \\
Initial biomass & $x_{4,\text{init}}$ & 60 g/m$^2$ & 8 g/m$^2$ & Est. \\
\bottomrule
\end{tabular}
\end{table}